\chapter{Interferenčné prúdy}

Interferenčná terapia (IFP), často označovaná ako "kráľovná elektroliečby", predstavuje jednu z najsofistikovanejších modalít vo fyzikálnej medicíne. Patrí do skupiny stredofrekvenčných prúdov, ktoré vďaka svojim unikátnym fyzikálnym vlastnostiam dokážu prekonať bariéru kože s minimálnym odporom a doručiť tak terapeuticky účinnú elektrickú energiu priamo do hlboko uložených tkanív – svalov, nervov a kĺbov. Táto kapitola ponúka hĺbkovú analýzu teoretických východísk, biofyzikálnych interakcií a klinických aplikácií interferenčných prúdov, pričom osobitný zreteľ je kladený na ich nezastupiteľné miesto v komplexnej rehabilitačnej starostlivosti o pacientov s reumatoidnou artritídou \parencite{Capko2010}.

\section{História a vývoj metódy}

Genéza tejto revolučnej metódy siaha do polovice 20. storočia a je neodmysliteľne spätá s menom rakúskeho fyzika a vynálezcu, Dr. Hansa Nemeca (1907–1981). V 40. rokoch minulého storočia čelila rehabilitácia vážnemu technickému limitu: vtedajšie nízkofrekvenčné prúdy (napr. faradický impulz alebo diadynamické prúdy) mali výborný liečebný efekt, ale narážali na vysokú impedanciu ľudskej kože.
Koža sa pri nízkych frekvenciách správa ako izolátor. Aby sa prúd dostal do hĺbky, bolo nutné použiť vysokú intenzitu, čo však vyvolávalo nepríjemné pálenie, bolesť a v krajných prípadoch až poleptanie pokožky pod elektródou. Pacienti túto liečbu často netolerovali \parencite{Capko2010}.

Dr. Nemec prišiel s geniálnym riešením, ktoré obišlo tento problém. Všimol si, že s rastúcou frekvenciou klesá odpor kože. Navrhol preto použiť dva nezávislé prúdové okruhy so strednou frekvenciou (okolo 4000 Hz), pre ktoré je koža takmer priehľadná. Kľúčovou myšlienkou bolo, že tieto dva prúdy sa "stretnú" až v hĺbke tkaniva. V mieste ich prekríženia dôjde k fyzikálnemu javu interferencie, čím vznikne nový, nízkofrekvenčný pulz priamo v cieľovom orgáne. Tento princíp "endogénneho generovania" liečivého signálu znamenal zásadný zlom. Kým prvé prístroje známe ako "Nemectron" využívali statickú interferenciu, neskorší technologický vývoj umožnil vznik dynamických vektorových polí, schopných ošetriť aj rozsiahle a anatomicky zložité oblasti \parencite{Podebradsky2009, Goats1990}.

\section{Fyzikálne princípy interferencie}

Fundamentálnym princípom metódy je superpozícia vlnení. Do organizmu sú privádzané dva sínusové striedavé prúdy strednej frekvencie pomocou štyroch elektród. Prvý okruh má konštantnú nosnú frekvenciu $f_1$ (napríklad 4000 Hz) a druhý okruh má frekvenciu $f_2$ mierne odlišnú (napríklad 4100 Hz).
Tieto prúdy prechádzajú tkanivom bez toho, aby pacienta dráždili. V mieste, kde sa ich dráhy pretnú (geometrický stred polohy elektród), dochádza k ich vzájomnému skladaniu – interferencii. Výsledkom je vznik amplitúdovej modulácie. Frekvencia tejto modulácie nie je náhodná; presne zodpovedá rozdielu frekvencií oboch nosných prúdov. Tento parameter nazývame \textbf{Amplitúdovo modulovaná frekvencia (AMF)}:
\[ AMF = |f_2 - f_1| = |4100 - 4000| = 100 \text{ Hz} \]
Z klinického hľadiska to znamená, že hoci tkanivom preteká prúd 4000 Hz (ktorý necítime), bunky v hĺbke reagujú na "rytmus" 100 Hz (ktorý lieči). Týmto spôsobom dokážeme doručiť účinok nízkej frekvencie (1–100 Hz) do hĺbky, kam by sa priama nízka frekvencia cez kožu nedostala \parencite{Robertson2006, Goats1990}.

\subsection{Fenomén prekonania kožného odporu}
Elektrické vlastnosti kože sú kľúčom k pochopeniu výhody IFP. Koža sa správa ako kapacitný odpor (kondenzátor). Jej impedancia ($Z$) je nepriamo úmerná frekvencii prúdu ($f$):
\[ Z \approx \frac{1}{f} \]
Pri klasickej frekvencii 50 Hz je odpor kože enormný, približne 3200 $\Omega$. To predstavuje bariéru, ktorá bráni prieniku energie.
Avšak pri použití "nemecovskej" frekvencie 4000 Hz tento odpor klesá strmo nadol, až na hodnotu okolo 40 $\Omega$. Prúd tak vstupuje do tela ľahko, "hladko", bez dráždenia voľných nervových zakončení v koži (nociceptorov). Pacient necíti žiadne pálenie povrchu, ale vníma príjemné, masívne vibrácie v hĺbke svalu alebo kĺbu. To umožňuje terapeutovi použiť podstatne vyššie intenzity prúdu než pri iných metódach, čo vedie k silnejšiemu biologickému účinku \parencite{Navratil2019, Watson2020}.

\section{Fyziologické účinky na organizmus}

Spektrum účinkov IFP je determinované predovšetkým voľbou frekvencie AMF. Reumatológia využíva najmä tieto tri hlavné efekty:

\subsection{Analgetický účinok (Tlmenie bolesti)}
Potlačenie bolesti je najčastejšou indikáciou. Mechanizmus je duálny a závisí od frekvencie:
\begin{enumerate}
    \item \textbf{Vrátková teória bolesti (Gate Control Theory):} Aplikácia vyšších frekvencií AMF (80 – 100 Hz) silne stimuluje hrubé myelinizované nervové vlákna typu A$\beta$, ktoré vedú dotykové a tlakové informácie. Tieto rýchle vzruchy prichádzajú do zadných rohov miechy skôr než signály bolesti vedené pomalými vláknami C. Tým sa aktivujú inhibičné interneuróny, ktoré "zatvoria vrátka" pre bolesť. Efekt je okamžitý, pacient cíti úľavu priamo počas procedúry \parencite{Kolar2012}.
    \item \textbf{Opioidný mechanizmus:} Pri nízkych, rytmických frekvenciách (1 – 10 Hz) dochádza k stimulácii produkcie endogénnych opiátov (endorfínov a enkefalínov) v centrálnom nervovom systéme. Tento typ analgézie nastupuje pomalšie (po cca 15-20 minútach), ale pretrváva dlhodobo, často niekoľko hodín po ukončení terapie \parencite{Navratil2019}.
\end{enumerate}

\subsection{Myorelaxačný účinok (Uvoľnenie svalov)}
Pri reumatoidnej artritíde sú svaly v okolí postihnutého kĺbu často v reflexnom kŕči (spazme), čo zvyšuje bolesť a obmedzuje pohyb. IFP využíva jav zvaný \textbf{Wedenského útlm}. Pri nosných frekvenciách okolo 4000 Hz nervové vlákno nie je schopné reagovať na každý jeden impulz depolarizáciou, pretože sa dostáva do stavu únavy (refraktérnej fázy). Dochádza k zníženiu dráždivosti a následnému uvoľneniu svalového tonusu \parencite{Podebradsky2009}.

\subsection{Vazodilatačný a antiedematózny účinok}
Zlepšenie prekrvenia je vitálne pre odvádzanie zápalových metabolitov (laktát, prostaglandíny, histamín). IFP tlmí aktivitu sympatického nervového systému, čo vedie k relaxácii hladkej svaloviny v stenách ciev a ich rozšíreniu (vazodilatácii). Súčasne, rytmické striedanie kontrakcie a relaxácie svalov pri nižších frekvenciách (10-50 Hz) funguje ako "svalová pumpa", ktorá mechanicky vytláča venóznu krv a lymfu z oblasti, čím efektívne redukuje opuchy (edémy) \parencite{Capko2010}.

\section{Metodika aplikácie a variácie}

Správna technika je podmienkou úspechu. V praxi rozlišujeme dva základné spôsoby aplikácie, ktoré sa líšia počtom elektród a miestom vzniku interferencie.

\subsection{Dvojpólová interferencia (Premodulovaná)}
Tento režim využíva len dve elektródy (jeden kanál). Špecifikom je, že k zmiešaniu prúdov a vzniku nízkofrekvenčnej obálky nedochádza v tele pacienta, ale už v samotnom prístroji. Z elektród teda vystupuje už hotový, modulovaný prúd.
\textbf{Výhody:} Je to ideálna technika pre ošetrenie malých plôch, trigger bodov alebo drobných kĺbov ruky, kam by sa štyri elektródy fyzicky nezmestili.
\textbf{Nevýhody:} Nakoľko prúd je modulovaný už pred vstupom do kože, čiastočne sa stráca výhoda ľahkého prieniku typická pre klasickú interferenciu. Hĺbkový efekt je preto o niečo slabší \parencite{Watson2020}.

\subsection{Štvorpólová interferencia (Klasická)}
Zlatý štandard pri liečbe veľkých kĺbov a chrbtice. Vyžaduje použitie štyroch elektród zapojených do dvoch krížiacich sa okruhov. Liečivý účinok vzniká presne v strede, v tzv. zóne interferencie.
Moderné prístroje umožňujú modifikovať tvar tohto poľa:
\begin{itemize}
    \item \textbf{Statické pole:} Zóna účinku je fixná, má tvar štvorlístka. Vhodné pre lokalizovanú bolesť.
    \item \textbf{Vektorové pole (Dynamické):} Prístroj cyklicky mení intenzitu jedného okruhu, čím dochádza k "rotácii" poľa. Interferenčná zóna sa tak pohybuje a premasíruje celú oblasť medzi elektródami. Toto je metóda voľby pri difúznych bolestiach kĺbov (napr. celé rameno, koleno) \parencite{Navratil2019}.
\end{itemize}

\subsection{Aplikácia s vákuovou jednotkou (Vákuová terapia)}
Veľmi často sa v klinickej praxi stretneme s kombináciou interferenčných prúdov a vákuovej jednotky. Elektródy nie sú pripevnené elastickými pásmi, ale sú umiestnené vo vnútri špeciálnych pohárikov (baniek), ktoré sa prisajú na kožu pomocou podtlaku.
Táto kombinácia prináša synergický efekt dvoch terapií v jednom čase:
\begin{enumerate}
    \item \textbf{Zlepšenie kontaktu:} Podtlak zabezpečí dokonalé priľnutie elektródy na nerovný povrch tela (napr. rameno, lopatka), čo zaručuje konštantný prenos prúdu.
    \item \textbf{Hyperémia (prekrvenie):} Samotný podtlak vtiahne kožu a podkožie do pohárika, čím spôsobí mechanickú vazodilatáciu. Zlepšené prekrvenie pod elektródou znižuje kožný odpor ešte výraznejšie než samotná frekvencia 4000 Hz, čím sa zvyšuje účinnosť prenosu energie do hĺbky.
    \item \textbf{Mikromasáž:} Moderné vákuové jednotky pracujú v pulznom režime, kedy sa sila prisatia rytmicky mení. Táto mechanická pulsácia pôsobí ako jemná vákuová masáž, ktorá stimuluje lymfatický obeh a pomáha pri redukcii opuchov.
\end{enumerate}
Pri použití vákuovej jednotky je však nutné dbať na opatrnosť u pacientov s krehkými cievami (varixy) alebo u seniorov s atrofickou kožou, kde by silný podtlak mohol spôsobiť petéchie alebo hematómy \parencite{Podebradsky2009}.

\section{Indikácie a kontraindikácie pri RA}

\subsection{Indikácie}
IFP je suverénnou metódou voľby najmä v subakútnom a chronickom štádiu RA. Hlavnými cieľmi sú:
\begin{enumerate}
    \item \textbf{Analgesia:} Zvládnutie bolesti bez nutnosti zvyšovať dávky liekov.
    \item \textbf{Resorpcia:} Urýchlenie vstrebávania chronických burzitíd a synoviálnych výpotkov.
    \item \textbf{Relaxácia:} Uvoľnenie reflexných spazmov paravertebrálneho svalstva a svalov v okolí postihnutých kĺbov.
    \item \textbf{Trofika:} Zlepšenie výživy tkanív pri akrocyanóze a Raynaudovom syndróme \parencite{Kolar2012, Capko2010}.
\end{enumerate}

\subsection{Kontraindikácie}
Hoci je IFP bezpečná, existujú stavy vylučujúce jej použitie:
\begin{itemize}
    \item \textbf{Absolútne:} Gravidita (aplikácia na driek a panvu), kardiostimulátor (riziko elektromagnetickej interferencie), febrilné stavy, akútna tromboflebitída a flebotrombóza (riziko embolizácie!), onkologické ochorenia v mieste aplikácie.
    \item \textbf{Relatívne/Lokálne:} Aplikácia v oblasti očí, karotických sínusov. Poruchy citlivosti kože vyžadujú opatrnosť. Prítomnosť kovových implantátov (TEP) nie je absolútnou kontraindikáciou, nakoľko stredofrekvenčný prúd nemá galvanickú zložku a nespôsobuje chemické popálenie, avšak odporúča sa nižšia intenzita kvôli riziku prehriatia kovu \parencite{Watson2020}.
\end{itemize}

\section{Porovnanie IFP s inými nízkofrekvenčnými prúdmi}
Pre správnu indikáciu je kľúčové chápať, v čom sa interferenčné prúdy líšia od klasických nízkofrekvenčných metód, ako sú TENS (Transkutánna elektrická nervová stimulácia) alebo diadynamické prúdy (DD).

\subsection{Interferencia vs. TENS}
TENS je najčastejšie používanou metódou domácej liečby bolesti. Pracuje s nízkou frekvenciou (1-150 Hz) priamo na koži.
\begin{itemize}
    \item \textbf{TENS:} Má silný povrchový účinok, ale ťažšie preniká do hĺbky kĺbu kvôli odporu kože. Pacient cíti ostré "bodanie" alebo mravčenie priamo pod elektródou.
    \item \textbf{IFP:} Vďaka nosnej frekvencii 4000 Hz prekonáva kožný odpor ľahšie. Účinok vzniká až v hĺbke ("endogénne"). Pacienti tolerujú vyššie intenzity prúdu pri IFP než pri TENS, čo vedie k hlbšiemu myorelaxačnému účinku \parencite{Goats1990}.
\end{itemize}

\subsection{Interferencia vs. Diadynamické prúdy (DD)}
DD prúdy (podľa Bernarda) sú staršou metódou využívajúcou sínusové prúdy 50/100 Hz.
\begin{itemize}
    \item \textbf{Tolerancia:} DD prúdy majú silnú galvanickú zložku (jednosmernú), ktorá dráždi kožu a spôsobuje hyperémiu (začervenanie) pod elektródou. Hrozí riziko poleptania kože pri dlhšej aplikácii.
    \item \textbf{Komfort:} IFP nemajú galvanickú zložku (sú čisto striedavé), preto nespôsobujú chemické zmeny pod elektródou ani pálenie. Sú bezpečnejšie pre pacientov s citlivou pokožkou a môžu sa aplikovať aj v prítomnosti kovových implantátov (s opatrnosťou), kým DD sú pri kovoch prísne kontraindikované \parencite{Podebradsky2009}.
\end{itemize}

\section{Súčasný stav dôkazov (Evidence-Based Medicine)}

V ére medicíny založenej na dôkazoch je nevyhnutné overovať empirické skúsenosti exaktnými dátami z klinických štúdií. Výskum účinnosti interferenčných prúdov v posledných rokoch priniesol viacero významných zistení, ktoré potvrdzujú ich miesto v moderných liečebných protokoloch.

\subsection{Účinnosť pri bolestiach chrbtice a kĺbov}
Najnovší systematický prehľad a meta-analýza, ktorú publikovali \textcite{Rampazo2023}, sa zamerala na pacientov s chronickou nešpecifickou bolesťou dolnej časti chrbta (LBP), čo je častá komorbidita u pacientov s reumatickými ochoreniami. Autori analyzovali dáta z viacerých randomizovaných kontrolovaných štúdií a dospeli k záveru, že IFP má štatisticky významný efekt na zníženie intenzity bolesti oproti placebu. Dôležitým zistením bolo, že IFP nielenže znižuje bolesť bezprostredne po aplikácii, ale zlepšuje aj funkčné parametre (pohyblivosť), čo je pre rehabilitáciu kľúčové.

Podobne povzbudivé výsledky priniesla štúdia zameraná na osteoartrózu kolena, ktorú viedli \textcite{Chen2022}. Hoci sa zamerali na gonartrózu, patofyziologický mechanizmus bolesti je podobný ako pri RA. Ich meta-analýza potvrdila, že interferenčná terapia je efektívna nielen v manažmente bolesti, ale významne prispieva k zlepšeniu kvality života pacientov. V porovnaní s inými elektroterapeutickými modalitami sa IFP javí ako superiorná najmä vďaka schopnosti preniknúť do hĺbky kĺbneho puzdra, čo povrchové prúdy (napr. TENS) nedokážu v takej miere.

\subsection{Neuromodulačný potenciál}
Zaujímavý pohľad prináša práca \textcite{Moore2018}, ktorá skúmala využitie IFP pri neuromodulácii gastrointestinálnych porúch. Hoci to priamo nesúvisí s liečbou kĺbov, potvrdzuje to schopnosť interferenčných prúdov ovplyvňovať autonómny nervový systém. Tento poznatok je aplikovateľný aj v reumatológii, kde modulácia sympatika (vazodilatačný efekt) zohráva úlohu pri liečbe trofických zmien a zlepšení mikrocirkulácie v zapálenom tkanive.

Všeobecne možno konštatovať, že IFP predstavuje bezpečnú, neinvazívnu alternatívu k farmakologickej liečbe bolesti, čím umožňuje znížiť spotrebu analgetík a nesteroidných antireumatík, ktoré majú pri dlhodobom užívaní vážne vedľajšie účinky \parencite{Goats1990}.

\section{Praktická aplikácia pri RA}

Aplikácia interferenčných prúdov u pacientov s reumatoidnou artritídou vyžaduje zohľadnenie špecifík tohto ochorenia. Na rozdiel od degeneratívnych artropatií, kde je zápalová zložka sekundárna, pri RA dominuje primárny imunitne mediovaný zápal. Preto je voľba parametrov stimulácie a načasovanie terapie kľúčová pre dosiahnutie optimálneho terapeutického efektu.

\subsection{Parametre stimulácie}
\begin{enumerate}
    \item \textbf{Akútna bolesť a edém:}
    Volíme vyššie frekvencie AMF, typicky 90 – 100 Hz, prípadne s malým rozkmitom (Spectrum) napr. 10 Hz. Cieľom je analgézia cez vrátkovú teóriu. Intenzita musí byť striktne podprahová (jemné mravčenie), nesmie vyvolať motorickú reakciu. Doba aplikácie kratšia, 10–15 minút.
    \item \textbf{Chronická stuhnutosť svalov:}
    Volíme nízke frekvencie AMF, 1 – 10 Hz, alebo striedavé režimy (1–100 Hz). Cieľom je myorelaxácia a "gymnastika ciev". Intenzita môže byť vyššia, až na hranicu motorického prahu (viditeľné chvenie svalu). Doba aplikácie dlhšia, 15–20 minút.
    \item \textbf{Zlepšenie trofiky tkanív:}
    Pri vaskulárnych komplikáciách RA (Raynaudov fenomén, akrocyanóza) volíme strednú frekvenciu 50 Hz s vektorovým poľom. Kombinujeme s podtlakovými elektródami pre synergický vazodilatačný efekt \parencite{Capko2010}.
\end{enumerate}

\subsection{Aplikačné schémy pre jednotlivé kĺby}

\textbf{Ramenný kĺb:}
Využíva sa tetrapolárna aplikácia s vektorovým poľom. Jedna elektróda sa umiestňuje ventrálne, druhá na úpon deltoidu, tretia a štvrtá dorzálne. Interferenčné pole tak prestúpi celým kĺbnym puzdrom. Ideálne je použitie vákuových elektród pre lepší kontakt na oblých tvaroch ramena. Pri frozen shoulder (adhezívna kapsulitída), ktorá môže komplikovať RA, volíme kombinovaný protokol: najprv analgézia (100 Hz, 10 min), potom myorelaxácia (1-10 Hz, 10 min) \parencite{Navratil2019}.

\textbf{Kolenný kĺb:}
Klasická aplikácia "do kríža". Elektródy obklopujú patelu. Pri oslabení kvadricepsu volíme frekvencie okolo 50 Hz, ktoré majú tonizačný vplyv na svalstvo. Pri opuchoch volíme 100 Hz na podporu rezorpcie. U pacientov s RA je častá prítomnosť Bakerovej cysty v popliteálnej oblasti; v takomto prípade je vhodné zahrnúť túto oblasť do terapeutického poľa \parencite{Kolar2012}.

\textbf{Ruky a zápästia:}
Pre ošetrenie drobných kĺbov rúk, ktoré sú pri RA najčastejšie postihnuté, je suverénne najlepšou metódou \textbf{"vodný kúpeľ"}. Do dvoch plastových nádob s vodou sa vložia uhlíkové elektródy. Pacient ponorí ruky (alebo nohy) do vody. Voda zabezpečí dokonalý kontakt s celým nerovným povrchom prstov. Táto metóda umožňuje ošetriť všetky kĺby ruky naraz, čo je pri lepení elektród nemožné \parencite{Navratil2019}.

\textbf{Bedrový kĺb:}
Hlboká lokalizácia bedrového kĺbu vyžaduje použitie väčších elektród a vyššej intenzity. Tetrapolárna aplikácia s elektródami umiestnenými v oblasti triesla, laterálne od veľkého trochanteru, gluteálne a posteromediálne. Vektorové pole zabezpečí ošetrenie celého kĺbu vrátane jeho hlbokých štruktúr \parencite{Robertson2006}.

\textbf{Krčná chrbtica:}
U pacientov s RA je nutná opatrnosť kvôli možnej atlantoaxiálnej instabilite. Pred aplikáciou IFP na krčnú chrbticu by mal byť vylúčený subluxácia C1-C2 pomocou RTG vyšetrenia. Aplikácia sa realizuje paravertebrálne s nízkymi intenzitami a frekvenciami 80-100 Hz pre analgéziu \parencite{Podebradsky2009}.

\subsection{Frekvencia a dĺžka liečby}
Štandardný liečebný cyklus pozostáva z 10-15 sedení aplikovaných denne alebo obdeň. Účinok jednej procedúry trvá niekoľko hodín, pri sériovej aplikácii dochádza ku kumulácii efektu. Po ukončení cyklu nasleduje prestávka 4-6 týždňov pred prípadným opakovaním. U pacientov s chronickou RA možno liečbu opakovať podľa potreby, typicky 2-3 cykly ročne \parencite{Capko2010}.

\subsection{Kombinácia s inými terapeutickými modalitami}
Interferenčné prúdy je možné efektívne kombinovať s ďalšími fyzioterapeutickými metódami. Najčastejšie sa IFP aplikuje pred kinezioterapiou – zníženie bolesti a uvoľnenie svalov umožňuje pacientovi efektívnejšie cvičiť. Kombinácia s hydroterapiou (teplý bazén) je obzvlášť výhodná – vztlak vody odľahčuje kĺby a teplo potencuje vazodilatačný efekt elektrického prúdu \parencite{guth2010rehabilitacia}.

Pri lokálnej aplikácii kryoterapie pred IFP u pacientov s akútnym vzplanutím využívame protizápalový efekt chladu nasledovaný analgetickým účinkom elektrického prúdu. Táto kombinácia je preferovaná pred termoterapiou v akútnej fáze zápalu \parencite{Podebradsky2009}.

\section{Bezpečnostné aspekty a nežiaduce účinky}

Interferenčná terapia patrí medzi najbezpečnejšie elektroterapeutické modality. Stredofrekvenčný charakter prúdu eliminuje riziko chemického popálenia pod elektródou, ktoré hrozí pri galvanických prúdoch. Napriek tomu je nevyhnutné dodržiavať bezpečnostné protokoly \parencite{Watson2020}.

\subsection{Možné nežiaduce účinky}
\begin{itemize}
    \item \textbf{Prechodná erytéma:} Začervenanie kože pod elektródami je fyziologickou reakciou na vazodilatáciu a nevyžaduje prerušenie liečby.
    \item \textbf{Svalová únava:} Pri nadmernej intenzite alebo príliš dlhej aplikácii s motorickou odpoveďou môže dôjsť k únave svalstva.
    \item \textbf{Nepríjemné pocity:} Pacienti môžu pociťovať mravčenie, vibrácie alebo tlak. Tieto pocity by nemali byť bolestivé – v takom prípade je nutné znížiť intenzitu.
    \item \textbf{Vertigo:} Pri aplikácii v oblasti krku sa vzácne môže objaviť závrať, čo vyžaduje okamžité ukončenie procedúry \parencite{Goats1990}.
\end{itemize}

\subsection{Monitorovanie pacienta}
Počas procedúry by mal terapeut pravidelne kontrolovať subjektívne pocity pacienta. Otázky by mali smerovať k intenzite vnímania prúdu, lokalizácii pocitov a komfortu. U pacientov s reumatoidnou artritídou je potrebné venovať pozornosť aj stavu kože – atrofická, ztenčená koža pri dlhodobej kortikoterapii je náchylnejšia na poškodenie \parencite{Watson2020}.

\section{Záver kapitoly}

Interferenčné prúdy predstavujú efektívnu a bezpečnú terapeutickú modalitu s pevným miestom v komplexnej rehabilitácii pacientov s reumatoidnou artritídou. Ich unikátna schopnosť doručiť terapeuticky účinnú nízkofrekvenčnú stimuláciu do hlboko uložených tkanív bez dráždenia povrchových štruktúr ich predurčuje pre liečbu kĺbov postihnutých zápalovým procesom.

Súčasná vedecká evidencia podporuje využitie IFP na zvládnutie bolesti, redukciu opuchov a prípravu pacienta na kinezioterapiu. V kombinácii s farmakologickou liečbou a komplexným rehabilitačným programom prispievajú k zlepšeniu funkčnej kapacity a kvality života pacientov s RA \parencite{Rampazo2022, Fuentes2010}.

Budúci výskum by mal priniesť štandardizované protokoly špecificky pre reumatoidnú artritídu a dlhodobé sledovanie efektivity v kontexte modernej biologickej liečby.
